\documentclass[10pt,a4paper]{amsart}
\usepackage[margin=19mm]{geometry}
\usepackage{amsmath,amssymb,mathtools}
\usepackage[scr=rsfs,cal=euler]{mathalfa}
\usepackage{xfrac}
\usepackage[unicode,hidelinks,bookmarks=false]{hyperref}
\newcommand{\ie}{i.e.\ }
\newcommand{\cf}{cf.\ }
\newcommand{\resp}{respectively\ }
\newcommand{\Subsection}{\subsection}
\providecommand{\nobreakdash}{\nobreak\hbox{-}}
\setlength{\emergencystretch}{2em}
\multlinegap0pt
\usepackage{mathtools}
\usepackage{xcolor}
\usepackage{amssymb}
\usepackage{enumerate}
\newtheorem{proposition}{Proposition}
\newcommand{\bb}{{\bf {b}}} % evidenziamo il campo b in grassetto
\newcommand{\hhh}{{\mathcal{H}}}
\title{Transport and flow for horizontal Sobolev contact velocities in Carnot groups}
\author{Gianluca Somma, Simone Verzellesi, Davide Vittone}
\date{Source: arXiv:2609.08738v1 (CC BY 4.0)}
\begin{document}
\maketitle

\noindent\emph{Source and licence.} Transport and flow for horizontal Sobolev contact velocities in Carnot groups. Version arXiv:2609.08738v1 (CC BY 4.0), distributed by arXiv under the Creative Commons Attribution 4.0 International licence, http://creativecommons.org/licenses/by/4.0/, as recorded in the arXiv licence metadata for this exact version and shown on the abstract page https://arxiv.org/abs/2609.08738v1. Full licence notice: \texttt{licenses/arxiv-2609.08738-CC-BY-4.0.txt}.\par\smallskip

\noindent\textit{Editorial scope.} Counted unit: the article's proposition prop\_contact\_components (source0.tex lines 555--561). The article's complete original proof of it is delivered with it (source0.tex lines 562--724, one \texttt{proof} environment of 163 lines). No mathematics is written, repaired or re-derived: the statement and the proof are the article's own lines, in the article's own notation, and no retained line is substituted or re-flowed.  Retained uncounted as the source-local context the proof needs: lines 220--222; lines 224--226; lines 252--254; lines 505--507; lines 517--519.  Nothing was omitted from any retained span, so the package carries no omission record.  No retained line contains a citation, so the wrapper carries no bibliography and no bibliography entry.  No retained line contains a figure environment, an image-inclusion command or drawing code, so no image is delivered and the package declares no asset dependency.  Editorial formatting only, in addition: an AMS \texttt{amsart} preamble that restates the article's own declarations and macros used by the retained lines, namely the environments \texttt{proposition} on separate counters, together with the standard packages those lines need.  The retained environments print in this document as Proposition 1 in order of appearance, whereas the article prints the same items with the numbering of its own full document.  The complete native source of the article as distributed in the arXiv e-print for this exact version is bundled unchanged as \texttt{sources/209742/source0.tex}.
\begin{equation}\label{defcarnotgroup}
\mathfrak{g}= V_1 \oplus \dots \oplus V_s, \qquad V_{i+1} = [V_1,V_i] \quad \text{for $i=1,\dots,s-1$,} \qquad [V_1,V_s]=\{0\}.
\end{equation}

\begin{equation} \label{eq_Jacobi}
[X,[Y,Z]]+[Y,[Z,X]]+[Z,[X,Y]]=0 \qquad \text{for every $X,Y,Z \in\mathfrak{g}$}
\end{equation}

\begin{equation} \label{eq_structural_constants}
\left[X^i_\alpha,X^j_\beta\right] \coloneqq \sum_{\gamma=1}^{n_{i+j}} c(i,j)^\gamma_{\alpha \beta} X^{i+j}_\gamma
\end{equation}

\begin{equation}\label{def_contact_SObolev_reg}
    \bb=\sum_{i=1}^s\sum_{\alpha=1}^{n_i}b^i_\alpha X^i_\alpha,\qquad b^i_\alpha\in W^{1,\theta}_{\hhh,\mathrm{loc}}(\mathbb G),\qquad [\bb,V_1]_p\subseteq\hhh_p\,\text{ for a.e.~$p\in\mathbb G$.}
\end{equation}

\begin{equation} \label{eq_contact_components_first}
X^1_\beta b^k_\gamma = \sum_{\alpha=1}^{n_{k-1}} c(k-1,1)^\gamma_{\alpha \beta} \, b^{k-1}_\alpha\qquad \text{for every $\beta=1,\dots,n_1$, $k=2,\dots,s$, $\gamma=1,\dots,n_k$.}
\end{equation}

\begin{proposition} \label{prop_contact_components}
  Let $\bb$ satisfy \eqref{def_contact_SObolev_reg}. Then
\begin{equation} \label{eq_contact_components}
X^j_\beta b^k_\gamma = \sum_{\alpha=1}^{n_{k-j}} c(k-j,j)^\gamma_{\alpha \beta} b^{k-j}_\alpha
\end{equation}
for every $j=1,\dots,s-1$, $\beta=1,\dots,n_j$, $k=j+1,\dots,s$ and $\gamma=1,\dots,n_k$.
\end{proposition}

\begin{proof}
We argue by induction on $j$. The case $j=1$ is true by \eqref{eq_contact_components_first}. Assume $j>1$, and assume that \eqref{eq_contact_components} holds for $j-1$ and for any $\beta=1,\dots,n_{j-1}$, $k=j,\dots,s$ and $\gamma=1,\dots,n_k$.
Fix $\delta=1,\dots,n_1$, $
\beta=1,\dots,n_{j-1}$, $
k=j+1,\dots,s$, 
$
\gamma=1,\dots,n_k$.
%Notice that 
%\begin{equation*}
%[X^1_\delta,X^{j-1}_\beta]b^k_\gamma
%=
%X^1_\delta\bigl(X^{j-1}_\beta b^k_\gamma\bigr)
%-
%X^{j-1}_\beta\bigl(X^1_\delta b^k_\gamma\bigr).
%\end{equation*}
By induction,
\begin{equation*}
X^{j-1}_\beta b^k_\gamma
=
\sum_{\alpha=1}^{n_{k-j+1}}
c(k-j+1,j-1)^\gamma_{\alpha\beta} \,
b^{k-j+1}_\alpha.
\end{equation*}
Therefore,
\begin{equation*}
X^1_\delta\left(X^{j-1}_\beta b^k_\gamma\right)
=
\sum_{\alpha=1}^{n_{k-j+1}}
c(k-j+1,j-1)^\gamma_{\alpha\beta} \,
X^1_\delta b^{k-j+1}_\alpha
\overset{\eqref{eq_contact_components_first}}{=}
\sum_{\alpha=1}^{n_{k-j+1}}
\sum_{\mu=1}^{n_{k-j}}
c(k-j+1,j-1)^\gamma_{\alpha\beta} \,
c(k-j,1)^\alpha_{\mu\delta} \,
b^{k-j}_\mu.
\end{equation*}
On the other hand, by \eqref{eq_contact_components_first},
\begin{equation*}
X^1_\delta b^k_\gamma
=
\sum_{\alpha=1}^{n_{k-1}}
c(k-1,1)^\gamma_{\alpha\delta} \,
b^{k-1}_\alpha.
\end{equation*}
Hence, by induction,
\begin{equation*}
X^{j-1}_\beta\left(X^1_\delta b^k_\gamma\right)
=
\sum_{\alpha=1}^{n_{k-1}}
c(k-1,1)^\gamma_{\alpha\delta} \,
X^{j-1}_\beta b^{k-1}_\alpha
=
\sum_{\alpha=1}^{n_{k-1}}
\sum_{\mu=1}^{n_{k-j}}
c(k-1,1)^\gamma_{\alpha\delta} \,
c(k-j,j-1)^\alpha_{\mu\beta} \,
b^{k-j}_\mu.
\end{equation*}
Combining the above equations,
\begin{equation*}
\left[X^1_\delta,X^{j-1}_\beta\right]b^k_\gamma
=
\sum_{\mu=1}^{n_{k-j}}
\left(
\sum_{\alpha=1}^{n_{k-j+1}}
c(k-j,1)^\alpha_{\mu\delta} \,
c(k-j+1,j-1)^\gamma_{\alpha\beta} \,
-
\sum_{\alpha=1}^{n_{k-1}}
c(k-j,j-1)^\alpha_{\mu\beta} \,
c(k-1,1)^\gamma_{\alpha\delta}
\right)b^{k-j}_\mu.
\end{equation*}
By the Jacobi identity \eqref{eq_Jacobi}, if $\mu=1,\ldots,n_{k-j}$,
\begin{equation*}
\left[\left[X^{k-j}_\mu,X^1_\delta\right],X^{j-1}_\beta\right]
-
\left[\left[X^{k-j}_\mu,X^{j-1}_\beta\right],X^1_\delta\right]
=
\left[X^{k-j}_\mu,\left[X^1_\delta,X^{j-1}_\beta\right]\right].
\end{equation*}
Applying \eqref{eq_structural_constants} to the above three terms, we obtain
\begin{align*}
\left[\left[X^{k-j}_\mu,X^1_\delta\right],X^{j-1}_\beta\right]
&=
\sum_{\alpha=1}^{n_{k-j+1}}
\sum_{\eta=1}^{n_k}
c(k-j,1)^\alpha_{\mu\delta} \,
c(k-j+1,j-1)^\eta_{\alpha\beta} \,
X^k_\eta,
\\
\left[\left[X^{k-j}_\mu,X^{j-1}_\beta\right],X^1_\delta\right]
&=
\sum_{\alpha=1}^{n_{k-1}}
\sum_{\eta=1}^{n_k}
c(k-j,j-1)^\alpha_{\mu\beta} \,
c(k-1,1)^\eta_{\alpha\delta} \,
X^k_\eta,
\\
\left[X^{k-j}_\mu,\left[X^1_\delta,X^{j-1}_\beta\right]\right]
&=
\sum_{\lambda=1}^{n_j}
\sum_{\eta=1}^{n_k}
c(1,j-1)^\lambda_{\delta\beta} \,
c(k-j,j)^\eta_{\mu\lambda} \,
X^k_\eta.
\end{align*}
Since $\left\{X^k_1,\dots,X^k_{n_k}\right\}$ is a basis for $V_k$, we conclude that 
\begin{equation*}
\sum_{\alpha=1}^{n_{k-j+1}}
c(k-j,1)^\alpha_{\mu\delta} \,
c(k-j+1,j-1)^\gamma_{\alpha\beta}
-
\sum_{\alpha=1}^{n_{k-1}}
c(k-j,j-1)^\alpha_{\mu\beta} \,
c(k-1,1)^\gamma_{\alpha\delta}
=
\sum_{\lambda=1}^{n_j}
c(1,j-1)^\lambda_{\delta\beta} \,
c(k-j,j)^\gamma_{\mu\lambda}.
\end{equation*}
In particular,
\begin{equation}\label{aux_improved_uno}
    \left[X^1_\delta,X^{j-1}_\beta\right]b^k_\gamma
=
\sum_{\lambda=1}^{n_j}
\sum_{\mu=1}^{n_{k-j}}
c(1,j-1)^\lambda_{\delta\beta} \,
c(k-j,j)^\gamma_{\mu\lambda} \,
b^{k-j}_\mu.
\end{equation}
On the other hand,
\begin{equation}\label{aux_improved_due}
\left[X^1_\delta,X^{j-1}_\beta\right]b^k_\gamma\overset{\eqref{eq_structural_constants}}{=}
\sum_{\lambda=1}^{n_j}
c(1,j-1)^\lambda_{\delta\beta} \, X^j_\lambda b^k_\gamma.
\end{equation}
Therefore, comparing \eqref{aux_improved_uno} and \eqref{aux_improved_due} and setting 
\begin{equation*}
    X=\sum_{\lambda=1}^{n_j}
\left(
X^j_\lambda b^k_\gamma
-
\sum_{\mu=1}^{n_{k-j}}
c(k-j,j)^\gamma_{\mu\lambda} \,
b^{k-j}_\mu
\right)X^j_\lambda,
\end{equation*}
we have 
\begin{equation*}
    \left\langle\left[X^1_\delta,X^{j-1}_\beta\right],X\right\rangle=\sum_{\lambda=1}^{n_j}
c(1,j-1)^\lambda_{\delta\beta}
\left(
X^j_\lambda b^k_\gamma
-
\sum_{\mu=1}^{n_{k-j}}
c(k-j,j)^\gamma_{\mu\lambda} \,
b^{k-j}_\mu
\right)=0
\end{equation*}
for every $\delta=1,\ldots,n_1$ and $\beta=1,\ldots,n_{j-1}$. This fact and \eqref{defcarnotgroup} grant that $X=0$, whence \eqref{eq_contact_components} follows.
\end{proof}

\end{document}
